\chapter{THEORETICAL BACKGROUND AND TECHNOLOGIES}

\section{Introduction}

Over the last decade, the web has shifted from a static document-centric approach to a dynamic, system based model. Of course, while the main focus of past models were centered around creating content, the real-time web's primary focus has been to synchronize states of clients in real-time. 

Most current web development frameworks depend upon a robust supporting infrastructure, including high bandwidth, and very low latency. 

The implementation of the proposed web application encompasses the development of the frontend, the backend services, and the database. The frontend involves the user interface components designed for high responsiveness. The backend handles the core business logic and event-driven architecture, while the database ensures persistent data storage. Furthermore, the underlying network infrastructure is of particular importance, as it is essential for establishing the reliable and seamless communication required in such demanding environments. 

To address the extreme nature of each scenario, the system was designed to be extremely fast and resilient. This chapter will explain the stacks of software used in building the system through the lens of the theoretical component. This chapter will describe the architecture of the system from the outside-in, beginning with the construction of the front-end, continuing with the back-end mechanisms, how it was deployed, and finally concluding with the security concerns of the system and then provide a historical view of the network.

\section{Web Architecture and Technology Stacks}

The structure as a whole will be based on a custom version of the MERN Stack (MongoDB, Express.js, React, Node.js) \cite{banker2011mongodb, express, freeman2019pro, tilkov2010node} with significant use of WebSockets in order to enable real time features. The reasons behind this selection of technologies are due to their built-in capabilities to perform non-blocking, asynchronous I/O operations which would allow for thousands of users to have simultaneous connections to the application at all times, while never blocking the main thread.

Table \ref{tab:core_tech} outlines an overview of the core technology stack used within each section of the infrastructure.

\begin{table}[htbp]
\centering
\caption{Core Technology Stacks utilized in the architecture}
\label{tab:core_tech}
\begin{tabular}{|l|p{10cm}|}
\hline
\textbf{Domain} & \textbf{Technology and Implementation Role} \\ \hline
Frontend Core & React 18 for component-based rendering and Vite for rapid bundling. \\ \hline
Styling & Tailwind CSS \cite{tailwind} for utility-based styling and executing the custom visual aesthetic. \\ \hline
State Management & Zustand \cite{zustand} for lightweight, decoupled global state management. \\ \hline
Backend Runtime & Node.js and Express.js \cite{express} for the asynchronous event loop and REST routing. \\ \hline
Database & MongoDB for denormalized, partition-tolerant persistent data storage. \\ \hline
Real-time Engine & Socket.io for persistent, full-duplex WebSocket communication between all nodes. \\ \hline
\end{tabular}
\end{table}

\section{Frontend Technologies and Communication Flow}

The new User Interface (UI) for such systems utilizes modern front-end technologies to help alleviate the cognitive load placed on personnel operating in stressful environments, adhering to established user interface design principles \cite{galitz2007}. To achieve this goal in the presentation layer, React 18 was selected to develop the UI as reusable components \cite{react2024docs}. Complementary, Vite was selected over other available solutions, as the development environment, due to its fast hot module replacement and optimized production build generation \cite{vite2024docs}. A utility-first styling system was implemented utilizing Tailwind CSS \cite{tailwind} to create a glassmorphic look with sufficient contrast and legibility while avoiding additional libraries that could slow the initial loading time of the application. Zustand \cite{zustand} was utilized for managing application state (primarily user sessions and authentication). It provides a simple to implement boilerplate to utilize and exposes an API that prevents unnecessary rendering of components subscribed to state changes. Greek and English localization were both supported through utilization of the i18next localization library \cite{i18next2024}. Bilingual staff members may require support for multiple languages in high-pressure operational environments. The most significant aspect of the client-side architecture is the manner in which the client and server communicate with one another. Most web-based applications utilize a ``pull'' communications model to allow for communication between the client and server. However, due to the fast-paced nature of transient event environments (e.g., kitchen's festival operations), the time delay associated with the ``pull'' model is too large for the client to wait for the requested data. To meet the need for a real-time connection, the Socket.io WebSocket protocol was implemented into the system. WebSockets provide a full-duplex, persistent TCP connection between the Node.js back end and the React front end \cite{fette2011websocket}. Once a WebSocket connection is established, the Node.js server and React front end can send data frames to each other with minimal overhead. This provides several advantages, including the ability for the server to push data updates (menus, orders, etc.) to the cashiers' displays without the clients having to request such data, thus reducing network traffic and latency.

\section{Backend Technology Stacks}

The backend system provides continuous access to stored data, has the capacity to handle a large number of users or requests with continued efficiency and reliability, and allows each component to operate independently should all communication be severed between the components' respective servers. A major advantage of this application was its development based on a philosophy of designing the architecture from an architectural standpoint as well as the theoretical perspective of using available memory effectively in high load environments and in synchronizing real-time data across multiple distributed computer systems in a very large computing environment. As such, if one is developing an enterprise-level scalable architecture for such applications, then a robust, fault tolerant architecture would most likely be required.

\subsection{The Concept of an Event}
This event-based architecture indicated the flow of data throughout the system did not occur based on a synchronized request/response sequence for each component that represented an individual part of the systems overall functionality, instead it generated events and reacted to events in the system's operational environment. In this context an event within the operational environment of a system represents a tangible or concrete mathematical operation which causes a change to the systems business logic state. For example, upon confirmation of a purchase (by selecting the ``Confirm Payment'' option) an event payload is generated representing a new order, including the current cart and timestamp. The backend then asynchronously notifies all relevant sub-systems and all active instances of the client at that moment of the change to the state of the system, functioning as a dispatcher. Additionally, decoupling of the functional components of the web application have enabled the creation of reactive dash-boards and notification managers.

\subsection{Node.js and MongoDB}
The back-end of our system was built using Node.js - an open source JavaScript runtime environment. Unlike most application servers, which create a new thread (or threads) per user on the server, Node.js utilizes only one thread and an event loop (reactor pattern). Therefore, while the server waits for an I/O-bound operation, such as inserting an order into the database, it can continue to process the next request of the cashier because it doesn’t block its own thread of execution. It can also insert the order into the database asynchronously (and in the background), and call the appropriate callback once the insert operation has been completed. We chose this architecture because many real-time order management systems are I/O-bound instead of CPU-bound. By utilizing a Node.js server, we can eliminate the need to create and maintain thousands of threads and can therefore provide a high quantity of concurrent web-socket connections - especially useful in environments with limited hardware, such as festival deployments.

For persistence of our data, we utilized MongoDB. Theoretically, the choice of a database for distributed systems is largely based upon the CAP theorem, which states that a data store can only provide two out of three of the following properties: consistency, availability, and partition-tolerance \cite{brewer2012cap}. Relational databases typically focus on maximizing consistency, so if the network goes down, they will lock writes to prevent divergent data. Because we have chosen to prioritize availability and partition-tolerance in our deployment, we decided to utilize MongoDB with configuration options to maximize availability. Additionally, we employed a denormalized document model. Instead of utilizing expensive relational joins to pull together all of the relevant details of an order from the database, we included all of the details of the order within a single document.

\section{Network Infrastructure}

The physical requirements for the use of our application necessitate that we take a ``local first'' approach to installing the application as opposed to other cloud-based applications. For instance, during festivals, there generally is no reliable Internet access available to users who are not located inside the festival grounds. In order to alleviate this problem, we have designed our application to be a self-contained application that is able to function locally without the need for an Internet connection. To make this feasible, we are utilizing the full MERN technology stack to deploy all of the necessary server infrastructure to dedicated servers in the area where the festival is being held. We will also utilize a local area network (intranet) to connect all of the tablet devices utilized by cashiers and cooks. We will accomplish this using a combination of both consumer-grade and enterprise-grade routers to create a dedicated local Wi-Fi network. This will enable us to eliminate the need for data to be routed over a mobile tower or the general public Internet backbone. Due to the fact that this local network will include the complete transactional flow of the application, it will be capable of avoiding any potential delays or latency that would result from the use of a cloud-based solution. As a further measure of redundancy, if we have very limited bandwidth over the Internet, we will utilize the secondary backup systems to periodically sync up with the primary systems.

\section{Application Security}

The security needs of the application architecture were taken into account as it was developed. Zero trust architecture was used for all web socket messages and all api calls made by the application are authenticating. The jwt stateless authentication approach has been used to achieve the goal of partition tolerance and to decrease the amount of session queries to the database. The server creates an access token with the user's information and the server creates an access token with the user's details and a digital signature when the user logs onto the site. Each subsequent request contains the token created when the user logged onto the site, and the token is stored on the user's computer, allowing the server to use a mathematical validation to determine if the token is valid or invalid, rather than the need to have the server verify the status of the token in a data storage system. The server is thereby relieved of additional workloads associated with verification, and provides the assurance that a user remains authenticated even if the server is unavailable as a result of planned server maintenance or unplanned server hardware failure.

\section{Comparative Analysis of Technology Stacks under High-Load}

To confirm that the MERN stack and the design decisions made with respect to architecture are theoretically correct, we will analyze how the behavior of web development stacks changes when under extremely high loads and also when there is extreme instability in the local area network (LAN) communications. The scenario to be used for this analysis is one where hundreds of users are attempting to synchronize all of their transactional data at the exact same moment.

In a traditional LAMP-based setup (Linux, Apache, MySQL, PHP), a new process is spawned for every PHP session. For each user's request, a completely new instance of the PHP interpreter is brought into memory. This results in the server having to use a great deal of RAM. Since servers only have a limited amount of RAM, this eventually leads to ``thrashing.'' Thrashing occurs when the operating system spends the majority of its time moving data out of memory to disk and vice versa. This renders the machine unusable.

Most large-scale business applications utilize the Java Spring framework which utilizes a thread-per-connection model. While threads consume significantly less memory than spawning processes, when you have many simultaneous connections (e.g. thousands of simultaneous connections) the CPU is spending the vast majority of its time switching contexts as opposed to actually executing the program. Therefore, threads do not provide a suitable solution for extremely high levels of I/O concurrency. In contrast, the MERN stack utilizes an asynchronous event loop, allowing the majority of the requests coming from the client to be processed by a single thread. As such, the MERN stack requires substantially less memory per connection, and does not incur the significant CPU overhead associated with the context switching of threads.

\section{The Evolution of Point of Sale (POS) Architectures and Related Work}

A description of the overall objective of the developed system requires a historical overview of transaction processing. In many ways, the evolution of POS technology has been driven by advancements in both connectivity and data accessibility. The first generation of POS systems (which existed from approximately 1970 through 1999) were essentially electronic cash registers. These systems were primarily stand-alone hardware devices that offered extremely high levels of reliability, however, they were essentially separate data islands that did not offer any real-time operational intelligence. The second generation of POS systems involved the introduction of client-server systems utilizing local area networks (LAN). While client-server architecture allowed for the centralization of data, creating LANs for temporary events proved to be logistically cumbersome, and the potential for introducing single points of failure became a significant issue. Cloud-based, native software currently represents the standard in the POS industry. With the exception of local maintenance issues, this model creates a serious problem in unstable environments. When the external network experiences problems, these systems ``freeze,'' thereby failing to meet one of the fundamental requirements of service availability. The emerging architecture models (one example being the model utilized in this thesis), which combine modern web technologies for both the interface and the logic, while deploying the server locally, provide a new generation of POS solutions. The hybrid approach enables the same advantages of rapid development using web frameworks as the application's logic and user interface are executed locally. 

Although there is an abundance of literature related to POS systems, the majority of the research is focused on either the use of analytics in retail environments or the enforcement of security regulations in a corporate environment. Additionally, although there have been numerous studies regarding the use of offline-first web architectures, the majority of the research lacks examples of how these architectures can be applied in a real-world scenario, especially when dealing with multi-user applications that require real-time synchronization in highly latencies conditions. The primary void within the existing body of knowledge and practical applications is the lack of a system that addresses the latency problem and utilizes modern, component-based web frameworks, along with an event-driven, local-first topology specifically designed to operate effectively despite the limitations presented by temporary, large-scale events. This thesis fills this void by providing a detailed description of a system that directly addresses this issue through the implementation and field-testing of the described architecture.